%==========================================================================

\begin{frame}[fragile]

  {\Huge Feature highlights}

  \vspace{10pt}

\end{frame}

%==========================================================================

\begin{frame}[fragile]{Self-Similar RangePolicy Interface}
   \textbf{Kokkos always allowed reuse of functions on host and device}

\begin{code}[keywords={awesome,template,TeamVectorRange,is_execution_space_v,class,const,KOKKOS_FUNCTION,KOKKOS_INLINE_FUNCTION,KOKKOS_LAMBDA,void,parallel_for,Kokkos,RangePolicy}]
  KOKKOS_INLINE_FUNCTION
  double awesome(double val, int a) {...}

  void bar_host() {
    double val = awesome(1.0, 2);
    parallel_for("bar", N, KOKKOS_LAMBDA(int i) {
      double val2 = awesome(val, 3);
    });
  }
\end{code}

\pause
\vspace{1cm}
\begin{center}
\textbf{But what if \texttt{awesome} contains parallelism itself?}
\end{center}
\end{frame}

\begin{frame}[fragile]{Self-Similar RangePolicy Interface}
   \textbf{You need to distinguish code paths!}
\begin{itemize}
  \item Different functions.
  \item Policy factory.
  \item \texttt{if constexpr}
\end{itemize}

\begin{code}[keywords={awesome,template,TeamVectorRange,is_execution_space_v,class,const,KOKKOS_FUNCTION,KOKKOS_INLINE_FUNCTION,KOKKOS_LAMBDA,void,parallel_for,Kokkos,RangePolicy}]
  template<class ExecOrTeam>
  KOKKOS_FUNCTION double awesome(const ExecOrTeam& exec, int M) {
    if constexpr (is_execution_space_v<ExecOrTeam>) {
      parallel_for(RangePolicy(exec, 0, M), KOKKOS_LAMBDA(int i) {
        ...
      });
    } else {
      parallel_for(TeamVectorRange(exec, 0, M), [&](int i) {
        ...
      });
  } }
\end{code}

\pause
\textbf{This still warns about calling host-only \texttt{parallel\_for} from host-device function!}
\end{frame}

\begin{frame}[fragile]{Self-Similar RangePolicy Interface}

  \textbf{Generalizes \texttt{RangePolicy} to support both execution policies and team-level ranges}

  \vspace{15pt}

  Enables writing generic code callable from both host and device:

  \begin{code}[keywords={template,class,const,KOKKOS_INLINE_FUNCTION,KOKKOS_LAMBDA,void,parallel_for,Kokkos,RangePolicy}]
    template <class Exec, class X, class Y>
    KOKKOS_INLINE_FUNCTION void sum_views(const Exec& exec, 
                                          const X& x, const Y& y) {
      Kokkos::parallel_for(Kokkos::RangePolicy(exec, 0, x.extent(0)),
        KOKKOS_LAMBDA(const int& i) { x(i) += y(i); });
    }
  \end{code}

  \vspace{10pt}

  \small
  \begin{itemize}
    \item Can be called with execution space on host: \texttt{sum\_views(exec\_space, x, y);}
    \item Can be called with team handle in device code: \texttt{sum\_views(team, x, y);}
  \end{itemize}

\end{frame}

%==========================================================================

\begin{frame}[fragile]{View with mdspan-style Template Arguments}


  \begin{itemize}
    \item \texttt{std::mdspan} uses different style of arguments than classic \texttt{View}
  \end{itemize}

  \begin{code}[keywords={DataType,LayoutType,MemorySpace,MemoryTraits,ElementType,ExtentsType,LayoutType,AccessorType}]
    View<DataType, LayoutType, MemorySpace, MemoryTraits>

    mdspan<ElementType, ExtentsType, LayoutType, AccessorType>
  \end{code}

  \begin{itemize}
     \item{\texttt{DataType} $\Rightarrow$ \texttt{ElementType} + \texttt{ExtentsType}}
     \item{\texttt{LayoutType} $\Rightarrow$ \texttt{LayoutType}}
     \item{\texttt{MemorySpace} + \texttt{MemoryTraits} $\Rightarrow$ \texttt{AccessorType}}
  \end{itemize}
  \vspace{15pt}

  \textbf{Kokkos Views now support mdspan-compatible template parameters}
  \begin{itemize}
    \item Enables specifying index types other than \texttt{size\_t} for potentially smaller View memory footprint
    \item Aligns with ISO C++ standard \texttt{std::mdspan}
  \end{itemize}

\end{frame}

\begin{frame}[fragile]{View with mdspan-style Template Arguments}
\textbf{Need to use \texttt{mdspan} compatible layout and accessors from Kokkos:}
\begin{code}[keywords={layout_left, layout_right, layout_stride, layout_left_padded, layout_right_padded, Accessor}]
   layout_left layout_right layout_stride
   Experimental::layout_right_padded<Pad> Experimental::layout_left_padded<Pad>
   Experimental::Accessor<ElementType, MemSpace, MemTraits>
\end{code}

\textbf{An Example:}
\begin{code}
  // before
  View<T**, LayoutLeft, HostSpace, MemoryTraits<>>
  // equivalent in new style:
  View<T, dextents<size_t, 2>, layout_left, Accessor<T, HostSpace, MemoryTraits<>>>
\end{code}

\textbf{Current Limitations:}
\begin{itemize}
  \item None of the template arguments are defaulted.
  \item Compiler warnings inside Kokkos when using \textit{signed} index types.
  \item Arbitrarily mixing static/dynamic extents is not well tested.
  \item Custom layouts and accessors not officially supported yet.
\end{itemize}
\end{frame}
%==========================================================================

\begin{frame}[fragile]{SFC64 PRNG Added to Random Algorithms}

  \textbf{New high-quality pseudo-random number generator with reproducible stream partitioning}

  \begin{itemize}
    \item $2^{64}$ streams from one 64-bit seed; period $\geq 2^{64}$ (expected $\sim 2^{255}$)
    \item Pool can be initialized \textit{in parallel}, faster than XorShift for large pools
    \item \texttt{seed\_offset} enables bit-reproducible distributed partitioning
    \item Drop-in replacement for \texttt{Random\_XorShift64\_Pool}
  \end{itemize}

  \begin{code}[keywords={Kokkos,Random_SFC64_Pool,parallel_for,rand,get_state,free_state,DefaultExecutionSpace,auto,uint64_t,double}]
    // Pool 0 handles streams [0, 500), pool 1 handles streams [500, 1000)
    // Identical results to a single pool of 1000 states with seed_offset 0.    
    Kokkos::Random_SFC64_Pool<> 
      pool0(seed, /*seed_offset=*/0,   /*num_states=*/500),
      pool1(seed, /*seed_offset=*/500, /*num_states=*/500);

    // ... use pool0 / pool1 exactly like Random_XorShift64_Pool
  \end{code}

\end{frame}

%==========================================================================

\begin{frame}[fragile]{Rank-1 MDRangePolicy Support}

  \textbf{Extends \texttt{MDRangePolicy} to support one-dimensional ranges}

  \vspace{15pt}

  \begin{itemize}
    \item Enables rank-1 construction: 
    
    \begin{code}[keywords={MDRangePolicy,ExecSpace,Rank}]
      MDRangePolicy<ExecSpace, Rank<1>>(0, N[, tile_size])
    \end{code}

    \item GPU backends (CUDA/HIP/SYCL) use tile size 256
    \item Host backends (OpenMP/Threads) use chunk-size logic like \texttt{RangePolicy}
    \item Includes full \texttt{parallel\_for} and \texttt{parallel\_reduce} support
  \end{itemize}

  \vspace{10pt}

\end{frame}

%==========================================================================

\begin{frame}[fragile]{create\_mirror\_view\_and\_copy Convenience Overload}

  \textbf{New single-argument overload for mirroring and copying Views}

  \vspace{15pt}

  \begin{code}[keywords={auto,create_mirror_view_and_copy}]
    // Old API (two arguments required)
    auto mirror = Kokkos::create_mirror_view_and_copy(space, view);
  \end{code}

  \begin{itemize}
    \item Problem: If \texttt{view} uses \texttt{SharedSpace} using \texttt{HostSpace} as \texttt{space} argument will allocate a new \texttt{View}
  \end{itemize}

\textbf{Always get a host accessible copy - avoiding unnecessary allocation:}
  \begin{code}[keywords={auto,create_mirror_view_and_copy}]
    // New convenience overload (one argument)
    auto mirror = Kokkos::create_mirror_view_and_copy(view);
  \end{code}

  \vspace{15pt}

  \small Consistent with \texttt{create\_mirror()} and \texttt{create\_mirror\_view()}, simplifying View management

\end{frame}

%==========================================================================

% Examples

% note: always keep the [fragile] for your frames!

%\begin{frame}[fragile]{Example list}
%  \begin{itemize}
%      \item Item 1
%      \item Item 2 with some \texttt{code}
%      \begin{itemize}
%        \item Sub-item 2.1
%        \item Sub-item 2.2
%      \end{itemize}
%  \end{itemize}
%\end{frame}

%\begin{frame}[fragile]{Example code}
%    \begin{code}[keywords={std}]
%        #include <iostream>
%
%        int main() {
%            std::cout << "hello world\n";
%        }
%    \end{code}
%\end{frame}

%\begin{frame}[fragile]{Example table}
%    \begin{center}
%        \begin{tabular}{l|l}
%            a & b \\\hline
%            c & d
%        \end{tabular}
%    \end{center}
%\end{frame}

%==========================================================================
